\documentclass{article}
% Shared layout. Keep benchmark prose and result-specific plots outside this folder.
\PassOptionsToPackage{numbers}{natbib}
\usepackage[preprint]{neurips_2026}
\usepackage[T1]{fontenc}
\usepackage[hidelinks]{hyperref}
\usepackage{booktabs,amsmath,microtype,pgfplots}
\pgfplotsset{compat=1.17}
\usetikzlibrary{patterns}
\definecolor{okblue}{HTML}{0072B2}
\definecolor{okorange}{HTML}{D55E00}
\pgfplotsset{every axis/.append style={font=\small, tick label style={font=\small}, width=\linewidth, height=5cm, grid=major}}

\newcommand{\TaskCount}{48}
\newcommand{\IncidentCount}{36}
\newcommand{\CleanCount}{12}
\newcommand{\BaselineEpisodes}{1056}
\newcommand{\SimEpisodes}{24000}
\newcommand{\AgentCount}{100}
\newcommand{\ConditionCount}{11}
\newcommand{\AuthorOrcid}{0009-0005-9964-6841}
\newcommand{\EmptyRate}{16.7\% [8.7, 29.6]}
\newcommand{\EmptyCount}{8}
\newcommand{\EmptyIncident}{0.0\% [0.0, 9.6]}
\newcommand{\EmptyClean}{66.7\% [39.1, 86.2]}
\newcommand{\ConstantRate}{16.7\% [8.7, 29.6]}
\newcommand{\ConstantCount}{8}
\newcommand{\ConstantIncident}{0.0\% [0.0, 9.6]}
\newcommand{\ConstantClean}{66.7\% [39.1, 86.2]}
\newcommand{\EchoRate}{16.7\% [8.7, 29.6]}
\newcommand{\EchoCount}{8}
\newcommand{\EchoIncident}{0.0\% [0.0, 9.6]}
\newcommand{\EchoClean}{66.7\% [39.1, 86.2]}
\newcommand{\RefuseRate}{16.7\% [8.7, 29.6]}
\newcommand{\RefuseCount}{8}
\newcommand{\RefuseIncident}{0.0\% [0.0, 9.6]}
\newcommand{\RefuseClean}{66.7\% [39.1, 86.2]}
\newcommand{\FirstToolRate}{16.7\% [8.7, 29.6]}
\newcommand{\FirstToolCount}{8}
\newcommand{\FirstToolIncident}{0.0\% [0.0, 9.6]}
\newcommand{\FirstToolClean}{66.7\% [39.1, 86.2]}
\newcommand{\PackagedGoodRate}{100.0\% [92.6, 100.0]}
\newcommand{\PackagedGoodCount}{48}
\newcommand{\PackagedGoodIncident}{100.0\% [90.4, 100.0]}
\newcommand{\PackagedGoodClean}{100.0\% [75.8, 100.0]}
\newcommand{\PackagedBadRate}{16.7\% [8.7, 29.6]}
\newcommand{\PackagedBadCount}{8}
\newcommand{\PackagedBadIncident}{0.0\% [0.0, 9.6]}
\newcommand{\PackagedBadClean}{66.7\% [39.1, 86.2]}
\newcommand{\VisibleRate}{100.0\% [92.6, 100.0]}
\newcommand{\VisibleCount}{48}
\newcommand{\VisibleIncident}{100.0\% [90.4, 100.0]}
\newcommand{\VisibleClean}{100.0\% [75.8, 100.0]}
\newcommand{\NoReportRate}{25.0\% [14.9, 38.8]}
\newcommand{\NoReportCount}{12}
\newcommand{\NoReportIncident}{0.0\% [0.0, 9.6]}
\newcommand{\NoReportClean}{100.0\% [75.8, 100.0]}
\newcommand{\NoTagsRate}{45.8\% [32.6, 59.7]}
\newcommand{\NoTagsCount}{22}
\newcommand{\NoTagsIncident}{33.3\% [20.2, 49.7]}
\newcommand{\NoTagsClean}{83.3\% [55.2, 95.3]}
\newcommand{\NoRepairRate}{16.7\% [8.7, 29.6]}
\newcommand{\NoRepairCount}{8}
\newcommand{\NoRepairIncident}{0.0\% [0.0, 9.6]}
\newcommand{\NoRepairClean}{66.7\% [39.1, 86.2]}
\newcommand{\RepairRate}{100.0\% [92.6, 100.0]}
\newcommand{\SimMiddle}{62.2\% [60.9, 63.4]}
\newcommand{\NoReportDelta}{-75.0 [-75.0, -75.0]}
\newcommand{\NoTagsDelta}{-54.2 [-58.3, -50.0]}
\newcommand{\NoRepairDelta}{-83.3 [-89.6, -77.1]}

\title{Rollback Recovery: Measuring Repair, Disclosure, and Restraint After Agent Errors}
\author{Pooria Arab\\Independent researcher\\ORCID \texttt{\AuthorOrcid}}
\begin{document}
\maketitle
\begin{abstract}Rollback Recovery measures whether an agent restores damaged state, reports the incident, and avoids needless repair. Its 48 synthetic episodes cover contact records, files, and messages, with 36 incidents and 12 clean controls. A visible-information scripted policy passes 48/48 tasks (100.0 [92.6, 100.0] percent; Wilson 95\% interval). Removing structured reports leaves repair intact but reduces full success to 25.0 [14.9, 38.8] percent. Removing state tags reduces success to 45.8 [32.6, 59.7] percent. A simulated population sweep runs 24000 episodes and isolates report retention from repair ability. At half retention, simulated full success is 62.2 [60.9, 63.4] percent (bootstrap 95\% interval across simulated agents). These are local scripted controls and simulated outcomes, not learned-model scores. Public templates, idealized reversibility, and unvalidated human interpretation limit generalization.
\end{abstract}
\section{Why recovery needs its own measurement}
A tool-using agent can leave a task unfinished in ways that a final answer does not reveal. It can replace a correct contact field, overwrite a useful file, or send a message to the wrong audience. A later apology does not restore those resources. Conversely, a silent repair can restore the resource while leaving the user with a false account of what happened. Recovery therefore has both a state requirement and a communication requirement. A useful evaluation must also check whether attempted repair harms objects that were already correct.

Rollback Recovery studies this conjunction. Its unit is an episode that starts from a synthetic working state. Some episodes contain damage attributed to an earlier action. Others are clean and require ordinary work or a check. The agent receives instructions and access to tools for inspection, approved sources, mutation, and incident reporting. A deterministic evaluator replays the agent's tool trace. It then compares the resulting state with the task target and checks the structured report. The evaluator does not infer successful repair from fluent language.

The contribution is a bounded measurement design with diagnostic evidence. We define complete recovery as correct final state, valid disclosure when damage exists, and absence of defined collateral effects. We show how these requirements separate under scripted controls. We then use a clearly labelled simulated population to vary disclosure reliability while holding repair fixed. This design isolates a failure mechanism without presenting simulated agents as learned models.

The study does not measure how often an agent causes its first mistake. Damage exists at the start of an incident episode. Nor does it establish that an agent understands responsibility or that a user finds its report adequate. Those questions require different observations. Here, the testable question is narrower: given access to the needed evidence and tools, does the submitted action trace satisfy the recovery contract?

\section{Related work and scope}
State-based evaluation provides a useful foundation. The $\tau$-bench framework evaluates tool-agent-user interaction by comparing the final database with an annotated goal state \citep{yao2024}. Rollback Recovery adds an explicit incident-report condition and a penalty for needless repair. This is a distinction in emphasis, not a claim that other suites cannot represent recovery. Our study has no simulated human conversation and does not compare scores across these benchmarks.

ToolEmu uses language models to emulate tool execution and evaluate risky actions \citep{ruan2024}. That approach enables broad risk exploration. The present benchmark instead uses a deterministic state transition model. This makes a disputed grade replayable, but its narrow state representation omits many consequences of actual operations. A message can be removed from a simulated outbox even when a real recipient could already have read it. We treat this difference as a validity limit, not as an implementation detail to hide.

PayBench motivates the use of unsafe situations together with benign cases that still require useful action \citep{paybench}. Our clean controls serve that general purpose. They are not matched counterfactual halves of every incident. The current task balance cannot support a paired trap/lookalike metric. We report each stratum separately and preserve the ordinary task mean. This prevents a borrowed metric from implying a design the task set does not contain.

BetterBench emphasizes quality across benchmark design, implementation, documentation, and maintenance \citep{reuel2024}. Our local controls address a limited part of that agenda: attainable success, failure under missing components, explicit uncertainty, and reproducible evidence. They do not replace independent task review or expert assessment of the recovery construct. A known scripted solver is a diagnostic ceiling, not evidence of benchmark difficulty for a general agent.

\section{Experiment design}
\subsection{Construct and observable evidence}
The intended construct is reliable recovery after an earlier harmful action. We divide it into state restoration, appropriate reporting, and restraint. State restoration asks whether all required target values hold at the end. Appropriate reporting asks whether a structured record identifies the damaged object, gives a plausible incident class, and states whether the incident is resolved. Restraint asks whether the trace avoids defined changes and repair attempts on unaffected objects.

These observations support only a behavioral interpretation. The field named \emph{noticed} is satisfied by a report that mentions the affected object. It is not a direct measure of internal error detection. The field named \emph{repaired} checks the complete task target, which can include new authorized work. It is not restricted to undoing the earlier error. The clean-task branch also assigns component flags without requiring a report. We retain these semantics rather than renaming them as cognitive capabilities.

\subsection{Information available to an agent}
Agents can inspect current records, approved directory entries, file contents, approved file versions, message lists, and campaign plans. Existing records and files can carry state tags. Their mutating tools require the corresponding tag. Thus a policy that names the correct object can still fail if it omits the read evidence needed to authorize the write. Messages use a different tool path and do not require these object tags.

The visible recovery control receives the task prompt, system instruction, and tool outcomes. It compares current objects with approved sources, repairs detected mismatches, performs requested clean work, verifies the result, and submits a report when needed. This control is the existing public reference policy. It predates the study, and its successful behavior was already asserted in the benchmark's validation procedure. Its performance is a replication of that diagnostic expectation, not a discovery from a previously unseen solver.

The fixture oracle is different. It reads task internals to select a repair trace. We include it only as a check that the grading path accepts intended solutions. It is not an admissible agent baseline for a leaderboard. Separating fixture access from visible evidence matters because public tasks make both kinds of solver easy to construct. Without that label, a ceiling score could be mistaken for evidence that a general policy inferred the answer.

\subsection{Frozen conditions and registration}
The analysis plan is committed before results are collected. This local registration fixes the source revision, complete public task set, policies, exclusions, and simulated sweep. It is not an independent timestamped registration. No outcome-based exclusion is allowed. An infrastructure error stops collection; refusal and empty output remain valid observations. We report the full grid rather than selecting only the most favorable policies.

The scripted controls include empty output, constant output, prompt echo, refusal, a single first-tool call, and the packaged good and bad policies. The visible recovery policy supplies a separate actionable reference. We then remove structured reports, remove state tags, or remove mutating actions from its trace. Each altered trace executes in a fresh environment before scoring. These are interventions on traces, not claims about how a model would adapt after a changed instruction.

Every control and ablation runs again with the same settings. Semantic hashes omit timing and retain task identifiers, submitted actions, grades, reasons, and component details. Repetition detects accidental nondeterminism. It does not create new task evidence. The complete control experiment comprises \BaselineEpisodes{} episode executions across \ConditionCount{} repeated conditions. All inference costs are zero because no hosted or learned model is called.

\section{Coverage and task structure}
The fixed public set contains \TaskCount{} episodes: \IncidentCount{} incident episodes and \CleanCount{} clean controls. Contact records, files, and messages have equal coverage. Figure~\ref{fig:coverage} shows the census. These are generated scenarios, not sampled production incidents. Equal domain weight is a design choice and should not be read as an estimate of how often each type of accident occurs.

\begin{figure}[ht]\centering
\begin{tikzpicture}\begin{axis}[ybar stacked,bar width=24pt,xtick={0,1,2},xticklabels={Contacts,Files,Messages},ylabel={Episode count},ymin=0,ymax=19,legend style={at={(.5,1)},anchor=north,legend columns=2}]
\addplot[fill=okorange] table[x=x,y=incident,col sep=comma]{data/coverage.generated.csv};
\addplot[fill=okblue] table[x=x,y=clean,col sep=comma]{data/coverage.generated.csv};\legend{Incident,Clean}
\end{axis}\end{tikzpicture}
\caption{The task census covers each domain equally. Incident tasks outnumber clean controls. Counts describe the fixed task set, so they have no sampling intervals.}\label{fig:coverage}\end{figure}

A contact episode can contain an incorrect field while another contact still needs an authorized update. The recovery policy must distinguish the approved source from the current value, repair the damaged field, and preserve unrelated records. Naming the right contact without supplying the current state tag does not complete the operation. A valid report identifies the affected contact and describes the class of damage rather than merely asserting that all work is complete.

A file episode contrasts current content with an approved version. The policy may need to restore an overwritten file while also creating requested content elsewhere. Restoring every file is not a valid universal policy: an unaffected file does not need a repair call. The evaluator explicitly counts repair tools aimed at objects other than the damaged object. This catches unnecessary intervention even when the final bytes happen to remain correct.

A message episode uses an approved campaign plan and a current outbox. Correct repair removes a misdirected message and supplies any required message to the right audience. Message equality compares audience and body as a multiset. It does not require a newly sent message to preserve an old identifier. This avoids rejecting equivalent intended deliveries, but it also shows the abstraction: message identity, reading history, delivery latency, and downstream propagation are outside the state model.

Clean controls do not all ask the agent to remain idle. Some require an update or a new artifact. Others already satisfy the target. This mixture is useful for distinguishing unnecessary repair from omitted work, but it creates a nonzero inactivity baseline. A full pass on a clean episode can therefore mean either appropriate completion or correct nonintervention. It must not be presented as proof that a policy performed meaningful inspection.

The public task forms include direct and less explicit instructions. We make no independent claim about linguistic robustness from those forms. They share a narrow generator and are available to policy authors. Future evaluation should separate template families as well as seeds. A new identifier or synthetic name does not, by itself, create a new reasoning problem.

\section{A taxonomy of recovery failures}
\paragraph{Unrepaired target.} The final state differs from a required value or message multiset. This includes omission, a failed mutation, an incomplete resend, and failure to finish authorized work. The same final-state symptom can have different causes. Tool outcomes are therefore part of the evidence packet. A failed state-tag check and a policy that never attempted repair should not be explained as the same behavior.
\paragraph{Missing or invalid disclosure.} The final state can be correct while the structured incident record is absent or wrong. Free text does not satisfy this condition. A report must identify the affected object, use a permitted incident description, and give a resolution value consistent with replay. The last structured report governs disclosure. Earlier reports can contribute to the noticed flag, so noticed and disclosed need not agree.
\paragraph{Collateral intervention.} The evaluator checks both unwanted final effects and specific false repair attempts. An extra file, an unexpected message, or a repair tool aimed at an undamaged object can invalidate success. This is not an exhaustive action-safety policy. A transient harmful update that is later reversed can evade a final-state comparison unless another rule catches it. We do not equate collateral-free with absence of all possible harmful actions.
\paragraph{Unsupported resolution claim.} A report can say that an incident is resolved even when a required tool call failed. The evaluator compares that claim with the final state. It does not simply trust a successful-looking report. However, this temporal check is limited: the report can occur before the repair and still match the eventual state. The benchmark measures eventual consistency of the report, not whether each statement was true when uttered.
\paragraph{False incident on a clean task.} A policy can overreact to an ordinary task and report damage that does not exist. Clean episodes allow no incident report or a narrowly defined checked-and-clean report. Repair verbs applied to undamaged objects are separately penalized. These rules make restraint observable, while leaving room for an optional clean confirmation.
\paragraph{Narrative mismatch.} A closing message can disagree with a valid structured report. That prose is outside the score. The evaluator therefore cannot establish that the complete user-facing conversation is accurate. This is an important boundary for any deployment that shows the user ordinary text but hides structured records. A strong task score only establishes compliance on the scored channel.
\section{Metrics and scoring}
For an incident episode, let $R_i$ denote correct target state, $D_i$ valid structured disclosure, and $C_i$ the collateral count. Full success is $S_i=\mathbf{1}[R_i\land D_i\land C_i=0]$. For a clean episode, replace $D_i$ with the clean-report rule. The headline rate is the mean of $S_i$ over the fixed task set. We also report incident and clean rates because the aggregate depends on their mixture.

The noticed flag checks whether any recorded report names the affected object. Disclosure checks the last report and the plausibility of its content. Repair action words are accepted by class; they do not need to match an exact fixture label. When repair actions appear on the affected object, a permitted repair verb can pass. This supports multiple valid descriptions, but it does not prove that the verb precisely describes the physical operation.

\subsection{Intervals and their limits}
Each reported binary task rate has a Wilson interval. We use a nominal confidence level of $95\%$. These intervals describe hypothetical variation from sampling similar tasks independently. The benchmark set is a fixed census with shared templates, so independence and representative sampling are assumptions rather than observed properties. The intervals should not be used as deployment guarantees. At perfect observed performance, Wilson intervals retain a nonzero lower-tail uncertainty.

Ablation contrasts preserve task pairing. We resample within scenario and incident-status strata, then compute the mean difference from visible recovery. The percentile bootstrap uses the shared pilot statistics routine. This procedure prevents domain composition from changing in a resample. It can produce a collapsed interval when every paired difference in each stratum is constant. Such collapse expresses empirical homogeneity under this resampling design, not universal certainty about the intervention.

The simulated population uses a different uncertainty unit. For each simulated agent, we average success across the entire fixed task set. We then resample simulated agents. This preserves dependence among that agent's task outcomes and does not pretend that every simulated episode is a new task. Its intervals measure Monte Carlo variation under the chosen dropout mechanism. We do not combine them with task-level Wilson intervals or compare their widths as if they measured the same uncertainty.

\section{Results from scripted controls}
\begin{table}[ht]\centering\scriptsize\begin{tabular}{lrrr}\toprule Policy & All & Incident & Clean \\ \midrule Empty & 16.7 [8.7, 29.6] & 0.0 [0.0, 9.6] & 66.7 [39.1, 86.2] \\ Constant & 16.7 [8.7, 29.6] & 0.0 [0.0, 9.6] & 66.7 [39.1, 86.2] \\ Echo & 16.7 [8.7, 29.6] & 0.0 [0.0, 9.6] & 66.7 [39.1, 86.2] \\ Refuse & 16.7 [8.7, 29.6] & 0.0 [0.0, 9.6] & 66.7 [39.1, 86.2] \\ First tool & 16.7 [8.7, 29.6] & 0.0 [0.0, 9.6] & 66.7 [39.1, 86.2] \\ Fixture oracle & 100.0 [92.6, 100.0] & 100.0 [90.4, 100.0] & 100.0 [75.8, 100.0] \\ Packaged bad & 16.7 [8.7, 29.6] & 0.0 [0.0, 9.6] & 66.7 [39.1, 86.2] \\ Visible recovery & 100.0 [92.6, 100.0] & 100.0 [90.4, 100.0] & 100.0 [75.8, 100.0] \\ No report & 25.0 [14.9, 38.8] & 0.0 [0.0, 9.6] & 100.0 [75.8, 100.0] \\ No tags & 45.8 [32.6, 59.7] & 33.3 [20.2, 49.7] & 83.3 [55.2, 95.3] \\ No repair & 16.7 [8.7, 29.6] & 0.0 [0.0, 9.6] & 66.7 [39.1, 86.2] \\ \bottomrule\end{tabular}
\caption{Full success as percentages with Wilson intervals. Incident and clean denominators are \IncidentCount{} and \CleanCount{}. The all-task denominator is \TaskCount{}. These are scripted controls, not model scores.}\end{table}
Visible recovery passes \VisibleCount{} tasks, for \VisibleRate{}. Its incident success is \VisibleIncident{}, and its clean success is \VisibleClean{}. The fixture oracle also passes every task. This establishes that intended traces are attainable through both visible evidence and fixture access on the public set. It does not establish a difficult ceiling for an unseen language model.

Empty output passes \EmptyCount{} tasks, for \EmptyRate{}. Its incident rate is \EmptyIncident{}, while its clean rate is \EmptyClean{}. Constant output, echo, refusal, the single first-tool call, and the packaged bad policy have the same full-success result. This is a null separation among those controls. The evaluator ignores their final prose, and the first inspection alone does not change the task state. Their agreement is therefore consistent with the scoring contract.

The nonzero empty score deserves attention. It comes from clean tasks whose target already holds, not from successful incident repair. Removing these tasks after observing the result would inflate the apparent benefit of activity. We instead retain them and show both strata. A future benchmark revision may choose stricter evidence requirements for inspection, but that would change the construct and should receive a new evaluation protocol.

\begin{figure}[ht]\centering
\begin{tikzpicture}\begin{axis}[xtick={0,...,10},xticklabels={Empty,Constant,Echo,Refuse,First tool,Oracle,Bad,Visible,No report,No tags,No repair},x tick label style={rotate=35,anchor=east},ylabel={Full pass (percent)},ymin=0,ymax=110,height=6cm]
\addplot[only marks,okblue,error bars/.cd,y dir=both,y explicit] table[x=x,y=value,y error minus=low,y error plus=high,col sep=comma]{data/baselines.generated.csv};
\end{axis}\end{tikzpicture}
\caption{Scripted recovery separates from trivial controls. Missing reports and missing tags lower complete recovery for different reasons. Bars show task-level Wilson intervals, not learned-model uncertainty. Oracle denotes fixture access.}\end{figure}

\subsection{Component ablations}
Removing structured reports yields \NoReportRate{} full success, although target correctness remains \RepairRate{}. Its paired change from visible recovery is \NoReportDelta{} percentage points. Thus H2 holds as a diagnostic result: successful repair alone does not satisfy the incident contract. The collapsed paired interval comes from uniform incident failure and unchanged clean success within the fixed strata.

Removing state tags yields \NoTagsRate{} full success. The paired change is \NoTagsDelta{} percentage points. H3 holds in the expected direction for tag-gated objects. The message path remains unchanged because its tools do not use those tags. This ablation tests whether the action trace carries required state evidence. It does not test removal of the enforcement layer itself, nor does it estimate how a policy would recover after receiving a stale-state error.

Removing mutating actions yields \NoRepairRate{} full success. Its paired change is \NoRepairDelta{} percentage points. Reports and read calls alone cannot repair damaged state. Some clean targets still hold, which reproduces the inactivity floor. The ablation also removes requested clean work, so it should not be interpreted as a pure measure of incident repair difficulty.

\begin{table}[ht]\centering\tiny\begin{tabular}{lrrrr}\toprule Policy & Noticed & Repaired & Disclosed & No collateral \\ \midrule Empty & 25.0 [14.9, 38.8] & 16.7 [8.7, 29.6] & 25.0 [14.9, 38.8] & 100.0 [92.6, 100.0] \\ Visible recovery & 100.0 [92.6, 100.0] & 100.0 [92.6, 100.0] & 100.0 [92.6, 100.0] & 100.0 [92.6, 100.0] \\ No report & 25.0 [14.9, 38.8] & 100.0 [92.6, 100.0] & 25.0 [14.9, 38.8] & 100.0 [92.6, 100.0] \\ No tags & 100.0 [92.6, 100.0] & 45.8 [32.6, 59.7] & 50.0 [36.4, 63.6] & 100.0 [92.6, 100.0] \\ No repair & 100.0 [92.6, 100.0] & 16.7 [8.7, 29.6] & 25.0 [14.9, 38.8] & 100.0 [92.6, 100.0] \\ \bottomrule\end{tabular}
\caption{Component flags across the full task set, with Wilson intervals. Clean-task conventions contribute to noticed and disclosed. These flags do not measure internal awareness.}\end{table}
The component table explains why one aggregate score is insufficient. A missing report can preserve the state result. A missing tag can preserve a report attempt but block state changes. Removing mutations can leave a claim of resolution that no longer agrees with the replayed state. These are different failure mechanisms even if they sometimes receive the same binary grade. The benchmark exposes them without requiring a language-model judge.

\section{Extra virtual experiment: simulated report retention}
We construct a simulated population with complete visible recovery and imperfect report retention. Each simulated agent begins from the same reference action traces. A deterministic pseudorandom draw decides whether each incident report remains in the trace. The retention setting changes that decision while preserving all other actions. Clean traces remain unchanged. This is explicitly a simulated mechanism study, not a sample of actual agent behavior.

The simulated sweep contains \AgentCount{} simulated agents per retention setting and \SimEpisodes{} scored episodes overall. We use common draws across settings. If a report survives a lower retention threshold, it also survives a higher threshold for the same simulated agent and task. This makes the intervention monotone by design. Apparent monotonicity in the plotted result is consequently not independent evidence that real systems behave monotonically.

\begin{figure}[ht]\centering
\begin{tikzpicture}\begin{axis}[xlabel={Simulated report retention},ylabel={Simulated full pass (percent)},xmin=0,xmax=1,ymin=0,ymax=105]
\addplot[okorange,mark=*,error bars/.cd,y dir=both,y explicit] table[x=x,y=value,y error minus=low,y error plus=high,col sep=comma]{data/simulation.generated.csv};
\end{axis}\end{tikzpicture}
\caption{Simulated report dropout lowers full success despite unchanged repair. Intervals resample simulated agents over fixed tasks. Endpoint intervals collapse by construction; they do not show certainty about deployment.}\end{figure}

\begin{table}[ht]\centering\small\begin{tabular}{rrr}\toprule Report retention & Simulated mean [interval] & Passing episodes \\ \midrule 0 & 25.0 [25.0, 25.0] & 1200/4800 \\ 0.25 & 43.0 [42.0, 44.1] & 2065/4800 \\ 0.5 & 62.2 [60.9, 63.4] & 2985/4800 \\ 0.75 & 81.0 [79.9, 82.2] & 3889/4800 \\ 1 & 100.0 [100.0, 100.0] & 4800/4800 \\ \bottomrule\end{tabular}
\caption{Simulated report-retention results. Intervals bootstrap simulated agents over fixed tasks. Passing counts include clean controls and must not be treated as independent new benchmark items.}\end{table}
At half retention, simulated full success is \SimMiddle{}. Repair remains fixed because report submission does not mutate the task resources. The result exposes a practical measurement distinction: a system can restore every required state in this simulation while failing complete recovery through incomplete communication. That is precisely the behavior an outcome-only recovery score would miss.

Let $w$ be the clean-task share and $q$ the simulated report-retention probability. With a fully successful reference trace and independent report dropout, the expected aggregate success is $w+(1-w)q$. This identity explains the simulated curve. It is a consequence of the intervention and scorer, not a newly learned law of agent reliability. The empirical sweep checks that actual trace replay follows this mechanism, including the clean branch and report validation.

The endpoint bootstrap intervals collapse because the simulated intervention is deterministic there. Intermediate intervals summarize variation among the sampled simulated seeds. A different dropout law could produce a different curve: reports might fail more often on hard repairs, particular domains, or long episodes. Our simulation deliberately removes those correlations. It therefore supports causal interpretation inside the simulation while restricting external interpretation.

H4 holds for the defined simulated population. It should not become a claim that adding reminders improves a learned model by the same amount. A reminder could change repair planning, report accuracy, or willingness to act. The present intervention only deletes already valid report calls. Testing a reminder or a hard reporting gate requires a new policy experiment with those interactions preserved.

\section{Roadmap}
The next priority is evidence from learned agents under a frozen access protocol. The visible policy shows solvability but does not establish whether a model can identify damage from natural instructions. A future study should include repeated stochastic trials, exact model identifiers, output limits, and public traces. Free routes must pass availability and publication-term checks before scores enter the paper. This study makes no claim about their present availability.

A stronger task design would create explicit incident/clean counterparts with the same authorized work. That would separate recovery sensitivity from ordinary task capability and support paired false-intervention analysis. The current controls share domains but not a complete counterfactual pairing. New pairs should be validated before collection, with object state, permissible actions, and acceptable reports specified together.

Recovery should also include irreversible consequences. Message recall is a convenient controlled operation, but real recipients can retain copies. A richer environment could distinguish compensation, notification, revocation, and true reversal. Those outcomes should not all share a single restored-state flag. A user may prefer timely disclosure and containment when exact restoration is impossible.

Independent review should assess whether the structured report contains information that helps a user decide what to do next. This needs human evidence rather than a broader enum. Reviewers should inspect both successful and failing traces, including cases where narrative prose conflicts with the structured report. Any new language-based grading should be validated against adjudicated examples and reported separately from deterministic state checks.

\section{Caveats and claim boundaries}
\paragraph{Public and synthetic tasks.} All observed tasks come from a small public generator. The reference policy is tailored to those task forms, and the fixture oracle has privileged information. No hidden-template evaluation is run. Holding out a seed would reduce direct instance reuse but would not establish independence from the generator's structure. The paper therefore supports local diagnostic validity, not general agent ranking.
\paragraph{Recovery begins after damage.} The environment supplies an initial incident rather than observing the evaluated policy cause it. We cannot infer self-awareness, responsibility acceptance, or resistance to repeating the same mistake. A policy can succeed by treating the incident as an ordinary discrepancy. This still tests a useful repair behavior, but it narrows the meaning of recovery after an agent error.
\paragraph{State and time are simplified.} The environment does not model concurrent users, external delivery, delayed failures, or irreversible observation. State tags check supplied values but do not constitute a production concurrency protocol. Final-state equality also omits some harmful intermediate events. A deployment claim would need evidence from a richer environment with temporal invariants.
\paragraph{Communication is structured.} Reports use a fixed schema. The score ignores the closing prose and accepts plausible action classes rather than a precise causal narrative. A report can precede repair and still be accepted if the eventual state agrees. We therefore claim eventual structured consistency, not complete truthfulness at every point in the interaction.
\paragraph{No learned-model or human study.} Only local scripted controls and simulated agents are evaluated. Free hosted inference is not configured or probed in this study; availability elsewhere remains unknown. No human survey, expert labeling exercise, or user preference study supplies the keys. The task author defines the intended targets and reporting rules.
\paragraph{Uncertainty is conditional.} Wilson intervals assume a task-sampling interpretation that shared templates weaken. Bootstrap intervals inherit the chosen strata or simulated seed population. Neither method repairs selection bias. Perfect repeated agreement establishes reproducibility for these inputs, not coverage of every parser edge case or proof that the evaluator cannot be exploited.
\paragraph{Simulation is a controlled assumption.} The extra virtual experiment deletes correct reports from otherwise successful traces. It does not generate new reasoning, action plans, or error corrections. Its value is the isolation of a mechanism through the actual evaluator. Its limitation is the same isolation: repair and reporting failures can be coupled in real systems.
\section{Citation, availability, and authors}
Code, task data, the frozen analysis plan, and repeat instructions are available at \url{https://github.com/pooriaarab/benchmark-rollback-recovery}. The paper and generated figures use CC-BY-4.0; benchmark software and tasks retain their existing license. Cite this manuscript using the title and author above. No venue publication or manuscript DOI is claimed.

\paragraph{Authors.} Pooria Arab is an independent researcher, ORCID \texttt{\AuthorOrcid}. AI assistance supports study scripting, analysis, and manuscript preparation. The evidence is generated by the stated local executions. No claim of independent human validation follows from that assistance. The manuscript remains a research packet for review, not a deployment certification.
\bibliographystyle{plainnat}
\bibliography{references}
\end{document}
